1: % Reference entries, U (AN-11).  Law 1: os/iword/core/codes.c iw_ed_attr
2: % (F8: [UND] ... [und], the cursor between them; F8 again moves past the off
3: % code), iw_attr_faces (underlined text drawn underlined), face_a.c (F8).
4: % Shots: ../wordperfect51-decomp goldens MINT01, mint MINT06.
5: % Undelete is P3 (bead initech-qvgd part 2, WL-0146): face_a.c
6: % IW_FK_undelete (F1 on the editing screen), edit() (the deletions kept:
7: % Backspace, Delete, a code deleted after its question, a block deleted;
8: % attribute codes are not kept, Reference p. 697; a run of deletions at one
9: % place is one deletion, any cursor movement starts a new one; three kept,
10: % the oldest dropped; 32K each, UNDO_CAP, a larger deletion refused with
11: % "ERROR: Document too large"), paint (the deletion shown at the cursor in
12: % reverse video, the document unchanged), undo_key (1 restores at the
13: % cursor, 2 or Up the previous deletion, Down the later one, F1 / Esc
14: % cancel).  Units of Measure: Setup records the units (spec/iword/
15: % settings.h); the screen goes on in inches (iw_fmt_inches).
16: % GATES: test-pubs-iword test-iword-codes test-iword-reveal-codes test-flair-iword-codes test-iword-undo test-flair-iword-undo test-iword-setup
17: \Gates{test-pubs-iword test-iword-codes test-iword-reveal-codes test-flair-iword-codes test-iword-undo test-flair-iword-undo test-iword-setup}
18: \begin{Feature}{Undelete}
19: Undelete puts back text you have deleted.  InitechWord keeps the last three
20: deletions, and shows each in turn at the cursor, so that you can choose the
21: one to restore.
22: 
23: To restore deleted text,
24: \begin{WPSteps}
25: \item Move the cursor to where the text is to go back.
26: \item Press \FK{Cancel}{F1} on the editing screen, with no menu or prompt
27: showing.  The status line shows \Scr{Undelete 1 Restore; 2 Previous
28: Deletion: 0}, and the last deletion appears at the cursor in reverse video.
29: \item Select \Ch{Restore}{1} to put that text back
30: \Or
31: Select \Ch{Previous Deletion}{2} to show the deletion before it, then
32: select \Ch{Restore}{1}.
33: \end{WPSteps}
34: \Documents{wpmenu}{F1,1=Restore}
35: \Documents{wpmenu}{F1,2=Previous Deletion}
36: The text is put back at the cursor, with the codes it held.
37: 
38: \BeginReplay{ref-undelete}
39: For example, type \TypeRun{Collate the reports by date and by department.}
40: Press \PressKey{Backspace} eleven times\Keys{Backspace Backspace Backspace Backspace Backspace Backspace Backspace Backspace Backspace Backspace}
41: to erase \textit{department.}, then press \PressKey{Home, Left} and
42: \PressKey{Delete} eight times\Keys{Delete Delete Delete Delete Delete Delete Delete}
43: to erase \textit{Collate }.  Press \PressKey{End} and then
44: \Press{Cancel}{F1}:
45: \WordScreen[rows=1]{ref-undelete-1}{The last deletion, shown at the cursor}
46: The deletion shown is the last one, \textit{Collate }.  Select
47: \Choose{Previous Deletion}{2}:
48: \WordScreen[rows=1]{ref-undelete-2}{The deletion before it}
49: Select \Choose{Restore}{1}: \textit{department.} is back at the end of the
50: line.
51: \EndReplay
52: 
53: \Notes
54: \Topic{Canceling}
55: Press \FK{Cancel}{F1} or \K{Esc} again to leave Undelete without restoring
56: anything.  The document is just as it was: the text shown in reverse video
57: was never in it.
58: 
59: \Topic{Codes}
60: Codes deleted with the text come back with it, except bold, underline and
61: the other appearance and size codes: those deleted with Backspace or Delete
62: are not kept.  A block deleted with Block and Delete, or with Move, comes
63: back whole, its codes included.
64: 
65: \Topic{Moving Through the Deletions}
66: At the Undelete prompt, \Ch{Previous Deletion}{2} and \K{Up} show the
67: deletion before the one shown, and \K{Down} the one after it.  Restore puts
68: back the one shown.
69: 
70: \Topic{Size of a Deletion}
71: A deletion holds up to 32K (32,768 bytes).  If you try to delete more at
72: once, as when deleting a very large block, the status line shows
73: \Scr{ERROR: Document too large} and nothing is deleted; delete the text in
74: smaller parts.
75: 
76: \Topic{What Is a Deletion}
77: Each character, code or block you delete becomes part of a deletion.
78: Characters deleted one after another at the same place with Backspace or
79: Delete make one deletion, so a word erased letter by letter comes back
80: whole.  Moving the cursor -- with the cursor keys or by searching -- ends a
81: deletion; the next thing you delete begins another.  When there are three
82: deletions and you delete more, the oldest is dropped.
83: 
84: \Topic{When Nothing Was Deleted}
85: If nothing has been deleted since InitechWord was started, Undelete shows
86: nothing at the cursor, and Restore does nothing.
87: 
88: \SeeAlso{Block, Delete; Cancel; Delete; Move}
89: \end{Feature}
90: 
91: \begin{Feature}{Underline}
92: Underline draws a line under text.  On the editing screen underlined text is
93: drawn underlined.
94: 
95: To underline text as you type it,
96: \begin{WPSteps}
97: \item Press \FK{Underline}{F8}.  InitechWord puts the codes \Code{[UND]} and
98: \Code{[und]} in the document, with the cursor between them.
99: \item Type the text.
100: \item Press \FK{Underline}{F8} again to move past \Code{[und]}.
101: \end{WPSteps}
102: \Documents{wpkey}{F8}
103: \begin{WPNote}To underline text you have already typed, block it and press
104: Underline (see \InRef{Block, Attributes}).\end{WPNote}
105: 
106: \BeginReplay{ref-underline}
107: For example, type \TypeRun{Please read the } (with the space), press
108: \Press{Underline}{F8}, type \TypeRun{attached}, press \Press{Underline}{F8}
109: and type \TypeRun{ memo.}:
110: \WordScreen[rows=1]{ref-underline}{An underlined word}
111: \EndReplay
112: 
113: \Notes
114: \Topic{Double Underline}
115: For a double line, choose Dbl Und (3) from the Appearance menu of the Font
116: key (see \InRef{Appearance}).
117: 
118: \Topic{Spaces}
119: On the editing screen, the underline is drawn under everything between the
120: two codes, spaces included.
121: 
122: \Topic{Underline From the Font Menu}
123: Undln (2) on the Appearance menu of the Font key does what F8 does.
124: 
125: \SeeAlso{Appearance; Block, Attributes; Bold; Codes}
126: \end{Feature}
127: 
128: \begin{Feature}{Units of Measure}
129: InitechWord measures in inches.  The status line gives Ln and Pos in inches
130: from the top and the left edge of the paper, and you type margins in inches
131: (see \InRef{Margins, Left and Right}).  Line spacing is given in lines (see
132: \InRef{Line Spacing}).
133: 
134: \Notes
135: \Topic{Characters and Lines}
136: In the standard font, Courier 10 Pitch, a character is a tenth of an inch
137: wide, and a single-spaced line is a sixth of an inch high: so Pos grows by
138: 0.1 inch with each character and Ln by 0.17 inch with each line.
139: 
140: \Topic{Inside the Document File}
141: In the document file, positions and margins are kept in WordPerfect units of
142: 1/1200 inch (see Appendix D).
143: 
144: \Topic{Rounding}
145: Ln and Pos are shown to two decimal places at most, without trailing zeros:
146: \Scr{1.5"}, \Scr{1.17"}.
147: 
148: \Topic{Setup}
149: Setup's Units of Measure screen records a unit for margins and tabs and one
150: for the status line: inches, centimeters, points, 1200ths of an inch or
151: WordPerfect 4.2 units (see \InRef{Setup}).  InitechWord keeps your choice
152: and shows it on that screen, but in this release the status line and the
153: Format screens show inches, and margins are typed in inches, whatever the
154: choice.
155: 
156: \SeeAlso{Margins, Left and Right; Setup; Status Line}
157: \end{Feature}
